\section{Introduction}

In this chapter we're going to start by determining the vectors
that can define the orbital plane.

\section{Direction of Angular Momentum}
Firstly we will start with the fact that the angular momentum is a constant.
This means the motion must be in a plane normal to the angular momentum. If the
motion wasn't within the plane then the direction of the angular momenum would
change, and so it would no longer be constant. This means the angular momentum
vector determines the normal to the plane.

\section{Direction of the Eccentricity Vector}

When the orbiting body is at periapsis the velocity is perpendicular to the position vector.
This is because the periapsis is defined as the point of closest approach to the central body.
If the velocity at that point were not perpendicular to the position vector, it would have a 
radial component either toward or away from the central mass.

A radial component toward the central mass would imply that the particle continues to approach the
central mass after this point, so it is not yet at periapsis.

A radial component away from the central mass would imply that the particle was closer to the focus
at an earlier time.

Both possibilities contradict the assumption that the point is the periapsis.
Hence the velocity at periapsis has no radial component and must be perpendicular to the position vector.